\documentclass[12pt]{article}
\usepackage[T2A]{fontenc}
\usepackage[utf8]{inputenc}
\usepackage[russian]{babel}
\usepackage{amsmath,amssymb,amsthm}
\usepackage{geometry}
\geometry{margin=2.5cm}
\usepackage{indentfirst}

\title{Квантовая коррекция ошибок на гиперболических кодах\\ с локальными проверками}
\author{Jane~Huang, Ли~Чэнь}
\date{2024}

\begin{document}
\maketitle

\section*{Аннотация}
Предложен класс квантовых LDPC-кодов на основе гиперболических поверхностей с ограниченной кривизной, допускающих локальные проверки стабилизаторов. Доказана теорема о пороге исправления ошибок, не зависящем от размера системы. Построен эффективный декодер на основе убеждений с тензорными сжатиями.

\vspace{0.5cm}
\noindent\textbf{Ключевые слова:} квантовые LDPC-коды; гиперболические поверхности; порог ошибок; стабилизаторы; декодирование.

\vspace{0.3cm}
\noindent\textbf{Раздел журнала:} Информатика, вычислительная техника и управление. \quad \textbf{Файл:} 1 из 2.

\section{Введение}
Квантовые LDPC-коды обеспечивают эффективное исправление ошибок при локальных проверках. Гиперболические поверхности дают естественную геометрию для таких кодов.

\section{Основной результат}
Порог исправления ошибок $p_c$ не зависит от размера системы $n$:
\begin{equation}
p_c = \text{const} > 0 \quad \text{при} \quad n \to \infty.
\end{equation}

\section{Заключение}
Обсуждается реализация кодов на сверхпроводящих схемах.

\end{document}